\scanpage{37}

% B01 | printed letterhead, upper left
% letterhead
\begin{center}
UNIVERSITÉ DE PARIS\\
---$\cdot$---\\
ÉCOLE NORMALE SUPÉRIEURE\\
---$\cdot$---\\
45, RUE D'ULM\\
ODÉON 06$\cdot$45\\
---$\cdot$---
\end{center}

% B02 | first line
Il résulte du lemme 1 que $\chi$ est surjective.

% B03 | first dash, homomorphism paragraph
--- Montrons que $\chi$ est \deleted{\unsure{surj.|aussi}} un homomorphisme : \inserted{Soient $g_1,g_2\in G$} Soit $\gamma_1$ un \deleted{la} chemin tel que $\partial_1(\gamma_1)=x_0$, $\partial_0(\gamma_1)=g_1x_0$, idem \unsure{p\textsuperscript{r}|pr} $\gamma_2$ tel que $\partial_1(\gamma_2)=x_0$, $\partial_0(\gamma_2)=g_2x_0$. On a $\partial_1(g_1\gamma_2)=g_1x_0$ et $\partial_0(g_1\gamma_2)=g_1g_2x_0$ : le chemin $\inserted{\gamma=}(\gamma_1,g_1\gamma_2)$ vérifie $\partial_1\gamma=x_0$, $\partial_0\gamma=g_1g_2x_0$. \deleted{\unsure{Or on a|On a}} On a $p\gamma=(p\gamma_1,pg_1\gamma_2)=(p\gamma_1,p\gamma_2)$ \deleted{\unsure{or|donc}} d'où $\chi(g_1g_2)=\chi(g_1)\chi(g_2)$

% B04 | second dash and fixed-vertex construction
--- Le noyau de $\chi$ contient évidemment le stabilisateur de $x_0$. Soit $x$ un sommet de $E$ et $s\in G$ tel que $sx=x$. Soit $\deleted{\unsure{\gamma}}\lambda$ un chemin de $E$ tel que $\partial_1(\deleted{\gamma}\lambda)=x_0$, $\partial_0\deleted{\gamma}\lambda=x$. \deleted{$\unsure{(\gamma,-s\gamma)|(\lambda,-s\lambda)}$} $\gamma=(\lambda,-s\lambda)$ est un chemin de $E$ tel que $\partial_1(\gamma)=x_0$, $\partial_0\gamma=sx_0$. d'où $\chi(s)=\operatorname{classe}(p(\lambda,-s\lambda))=\operatorname{classe}(p(\lambda,-\lambda))=e\,(\text{élément neutre})$.

% B05 | paragraph beginning Le noyau, lower half
Le noyau de $\chi$ contient donc le sousgroupe $S$ \deleted{\unsure{des}} engendré par les stabilisateurs des \deleted{\unsure{points}} sommets de $E$. Ce sous groupe est distingué puisque le conjugué du stabilisateur de $x$ est le stabilisateur d'un transformé de $x$.

% B06 | third dash
--- On sait que $\operatorname{Ker}\chi\supset S$. Montrons qu'il ne le contient pas strictement.

% B07 | Lemme 3 at foot
\uline{Lemme 3} : $E/S$ est un revêtement \inserted{normal} de $E/G$, de groupe $G/S$.

% archive sheet number: 36

